\section{Introduction}

Mangoes hold immense economic and cultural importance in the Philippines, serving
not only as the national fruit but also as a cornerstone of the agricultural sector.
As the third most essential fruit crop in the country, mango production reached an
estimated 596.34 thousand metric tons between April and June 2023, covering a
planted area of 184.03 thousand hectares \cite{philippineMangoIndustry}. To maintain
global competitiveness and ensure high-quality produce, strict quality assessment
and sorting mechanisms are required during post-harvest processing. These
processes traditionally depend on manual sorting based on weight, size, and ripeness,
which are highly laborious, time-consuming, and prone to human error. Developing a
non-destructive and efficient method for estimating the weight, volume, and other
morphological properties of mangoes before harvesting or during processing is
critical to improving agricultural efficiency and yield estimation.

\subsection{Problem Statement}
One of the primary challenges in the automated visual inspection of green fruits is
green-to-green segmentation---the difficulty of distinguishing the green fruit from
similarly colored leaves and branches in the background \cite{depthcl}. This issue is
particularly pronounced in the Indian mango (\textit{Katchamitha}) \cite{Katchamitha2025},
which often retains a deceptively green skin even when fully ripe, occasionally
showing only a slight yellow blush \cite{Katchamitha20252}. Consequently, relying
solely on traditional RGB color space for segmentation and ripeness estimation is
unreliable. 

Furthermore, yield estimation and quality grading are hindered by the lack of non-
destructive measurement techniques. For instance, determining the weight or the
stone-to-flesh ratio of a mango typically requires destructive interference or
physical harvesting. Existing vision-based weight estimation models for mangoes
largely rely on 2D images \cite{volume2020mango, shape2015mango}, which inherently
neglect the actual 3D shape and partial depth of the fruit, leading to inaccuracies.
Research specifically addressing non-destructive stone-to-flesh ratio estimation
remains notably scarce.

\subsection{Objectives}
This study aims to develop a comprehensive, dual-pipeline framework for the non-destructive estimation of the morphological properties of Katchamitha mangoes. Specifically, the study seeks to:
\begin{enumerate}
    \item Utilize Monocular Depth Estimation (MDE), specifically the metric version of
    Depth Anything V2 \cite{depth_anything_v2}, to assist and refine the RGB instance
    segmentation performed by YOLOv11n (RGB + Depth Fusion).
    \item Implement and compare two processing architectures: (a) an uncalibrated
    pipeline relying on raw depth features and scale-invariant ratios, and (b) a
    calibrated pipeline using ArUco reference markers to translate pixel depth maps
    into physical dimensions.
    \item Extract robust geometric and color features from the fused RGB-Depth data
    to train an orientation classification model, recognizing that varying viewing
    angles require dedicated morphology and ripeness estimation pipelines.
    \item Train a multiple multi-output Random Forest regression model that
    estimates the mango's morphological properties based on the RGB-Depth
    information, the classified orientation, and the ripening phase. The
    properties to be estimated include:
    \begin{itemize}
        \item Dimensions (Height, Width, and Thickness), Volume, Weight, and Density of the mango.
        \item Dimensions, Volume, Weight, and Density of the mango's bone (stone).
        \item The ripening phase of the mango (ripe vs. unripe).
        \item The stone-to-flesh ratio of the mango.
    \end{itemize}
\end{enumerate}

\subsection{Significance of the Study}
The integration of MDE into a computer vision pipeline provides a low-cost,
adaptable alternative to expensive 3D scanning or LiDAR systems for automated
agriculture \cite{depthanythingv2}. By precisely estimating internal and external
properties without physically harming the fruit, this research contributes to
sustainable precision agriculture, minimizes post-harvest waste, and supports
local farmers in maximizing the market value of their yield.

\subsection{Scope and Limitations}
The dataset is limited to 45 Katchamitha mango samples gathered from the local
region of Mabalacat City, Pampanga. These 45 samples yield a total of 243 raw
images across six orientations, expanded to 640 images after data augmentation.
The dataset is split into two groups: 35 samples (482 augmented images) are used
for the uncalibrated pipeline, and 10 samples (158 augmented images) are equipped
with ArUco reference markers for the calibrated pipeline. The Katchamitha variant
was selected primarily for its local accessibility and cost-efficiency; however,
the proposed methodology is designed to be easily replicated for more commercially
dominant variants like the Carabao mango \cite{fernandez2017mango}.

Additionally, the metric version of MDE models can yield inconsistent absolute depth
values for the same object under varying lighting conditions or camera distances.
To mitigate this in the uncalibrated pipeline, the study relies on scale-invariant
depth ratios rather than raw depth maps, while the calibrated pipeline utilizes
the ArUco marker as a physical reference to establish an absolute scale factor.
Finally, the proposed system processes a single, unoccluded mango instance at a
time and is not yet trained to handle occluded environments in complex orchard
canopies.
